\documentclass[12pt]{report} % ---- COMPILE WITH XeLaTeX for Unicode!!! ----
\usepackage[a4paper, top=1in, bottom=1in, left=1.5in, right=1in]{geometry} % Set margins
\usepackage{setspace} % allows \doublespacing & \singlespacing
\usepackage{parskip} % enables non-indented paragraphs with a blank line

% enables graphics, all figures and images should be in the /figures/ folder:
\usepackage{graphicx}
\graphicspath{{figures/}}

% enhanced maths symbols and syntax:
\usepackage{amsmath}
\usepackage{amssymb}
%\usepackage{amsthm} % if you need mathematical definitions and theorems

\usepackage[ruled]{algorithm2e} % allow for algorithms in pseudocode
\usepackage{listings} % allow for source code algorithms/listings
\lstset{basicstyle=\footnotesize, frame=single, numbers=left} % customise listing style

% Customise chapter headings so they are of the format: 1. Introduction
\usepackage{titlesec}
\titleformat{\chapter}[block]
  {\normalfont\huge\bfseries}{\thechapter.}{1em}{\Huge}
\titlespacing*{\chapter}{0pt}{-20pt}{20pt}

% Setup bibliography - add file to project named bibliography.bib
% can create this from Zotero, Mendeley, refworks, etc
\usepackage{natbib}
\bibliographystyle{agsm}
\setcitestyle{authoryear,open={(},close={)}, aysep={,}}

\usepackage{hyperref} % add links to document

\doublespacing
\begin{document}
\input{titlepage} % Custom title page, see titlepage.tex for details
\pagenumbering{roman} % Roman numeral page numbering up to the first page of introduction
\setcounter{page}{2} % Start at "ii", to count title page as "i"

\chapter*{Abstract} % \chapter* means the chapter is not in the contents or numbered
Abstract goes here.

\chapter*{Acknowledgements}
I want to thank...

\tableofcontents

\chapter{Introduction}
\pagenumbering{arabic} % Introduction starts at page 1.

Introduction goes here. To cite like this:

\cite{johnson_perceptual_2016} showed that...

Use the "\textbackslash cite\{\}" command

\hfill \break

To cite like this:

A recent paper \citep{johnson_perceptual_2016} has shown that...

use the "\textbackslash citep\{\} command

\hfill \break

For papers with exactly three authors, you must use \textbackslash citep*\{\}, like this:

The Deep Learning review article \citep*{lecun_deep_2015} summarised...


\chapter{Background}

For an example figure, see Figure \ref{fig:keelelogo}:
\begin{figure}[h]
\centering
\includegraphics[scale = 0.5]{keelelogo} %normally replace scale=
\caption{The Keele University logo.}
\label{fig:keelelogo}
\end{figure}

Example equation:
\begin{equation} \label{eq:1}
    \min_{G} \max_{D} \mathcal{L}(D, G) = \mathbb{E}_{\boldsymbol{x} \sim p_{data(\boldsymbol{x})}} \log D(\boldsymbol{x}) + \mathbb{E}_{\boldsymbol{z} \sim p_{\boldsymbol{z}(\boldsymbol{z})}} \log (1-D(G(\boldsymbol{z}))).
\end{equation}

Example pseudocode algorithm:

\begin{algorithm}[H] \label{alg:1}
\DontPrintSemicolon % remove this to add semicolons at the end of each line
 \For{number of training iterations}{
  Step 1\;
  Step 2\;
  Step 3\;
 }
 Step 4\;
 \caption{Training process for Algorithm 1...}
\end{algorithm}
\newpage
Example source code listing:
\begin{lstlisting}[language=Python, caption=Python example]
import numpy as np
 
def incmatrix(genl1,genl2):
    m = len(genl1)
    n = len(genl2)
    M = None #to become the incidence matrix
    VT = np.zeros((n*m,1), int)  #dummy variable
 
    #compute the bitwise xor matrix
    M1 = bitxormatrix(genl1)
    M2 = np.triu(bitxormatrix(genl2),1) 
 
    for i in range(m-1):
        for j in range(i+1, m):
            [r,c] = np.where(M2 == M1[i,j])
            for k in range(len(r)):
                VT[(i)*n + r[k]] = 1;
                VT[(i)*n + c[k]] = 1;
                VT[(j)*n + r[k]] = 1;
                VT[(j)*n + c[k]] = 1;
 
                if M is None:
                    M = np.copy(VT)
                else:
                    M = np.concatenate((M, VT), 1)
 
                VT = np.zeros((n*m,1), int)
 
    return M
\end{lstlisting}

Example of multiple equations:
\begin{align} \label{eq:2}
    \min_{D} \mathcal{L}_{LSGAN}(D) = & \frac{1}{2}\mathbb{E}_{\boldsymbol{x} \sim p_{data(\boldsymbol{x})}} (D(\boldsymbol{x}) - 1)^2 + \frac{1}{2}\mathbb{E}_{\boldsymbol{z} \sim p_{\boldsymbol{z}(\boldsymbol{z})}} (D(G(\boldsymbol{z})))^2
    \\
    \min_{G} \mathcal{L}_{LSGAN}(G) = & \frac{1}{2}\mathbb{E}_{\boldsymbol{z} \sim p_{\boldsymbol{z}(\boldsymbol{z})}} (D(G(\boldsymbol{z}))-1)^2.
\end{align}

Example table:
\begin{table}[h]
\begin{tabular}{l | c c c}
& \textbf{Per-Pixel} & \textbf{Per-Class} &\\
\textbf{Algorithm} & \textbf{Accuracy} & \textbf{Accuracy} & \textbf{Class IoU}\\\hline
Algorithm 2 (ours) & \textbf{0.53} & \textbf{0.16} & \textbf{0.11} \\
Algorithm 2 (original) & 0.42 & 0.12 & 0.08 \\
Baseline algorithm & 0.49 & 0.13 & 0.10 \\
\end{tabular}
\label{tab:tablelabel}
\caption{Results table for...}
\end{table}

\section{Section 1}

Insert text here...

\subsection{Subsection 1}

Insert text here...

\section{Section 2}

Insert text here...

\chapter{Main Body 1}

Insert text here...

\section{Section 1}

Insert text here...

\section{Section 2}

Insert text here...

\chapter{Evaluation}

Insert text here...

\section{Testing}

Insert text here...

\section{Experiments and Results}

Insert text here...

\section{Discussion}

Insert text here...

\chapter{Conclusion}

Insert text here...

\bibliography{bibliography.bib}

\appendix

\chapter{Project Plan}

\chapter{Ethics Form}

\end{document}